% !TEX root = ../reporte_final.tex
\section{Apéndice: código fuente esencial}
\label{ap:codigo}

Este apéndice contiene los fragmentos esenciales de Python que reproducen el experimento de principio a fin. El notebook completo se encuentra en \texttt{notebook.ipynb} dentro del repositorio del proyecto.

\subsection{Carga y preprocesamiento}

\begin{lstlisting}[language=Python,caption={Carga, limpieza y estandarización del dataset.}]
import os, urllib.request
import pandas as pd
from sklearn.preprocessing import StandardScaler, LabelEncoder

RANDOM_STATE = 42
NUM_COLS = ["bill_length_mm", "bill_depth_mm",
            "flipper_length_mm", "body_mass_g"]

DATA_URL = ("https://raw.githubusercontent.com/allisonhorst/"
            "palmerpenguins/main/inst/extdata/penguins.csv")
LOCAL_PATH = "data/penguins.csv"

if not os.path.exists(LOCAL_PATH):
    os.makedirs("data", exist_ok=True)
    urllib.request.urlretrieve(DATA_URL, LOCAL_PATH)

df = pd.read_csv(LOCAL_PATH)
df_clean = df.dropna(subset=NUM_COLS + ["species"]).reset_index(drop=True)

X = StandardScaler().fit_transform(df_clean[NUM_COLS])
y_true = LabelEncoder().fit_transform(df_clean["species"])
\end{lstlisting}

% Reservamos espacio para subsection + listing 3 completo, evitando
% que el titulo A.2 quede huerfano al pie de la pagina previa.
\Needspace{28\baselineskip}
\subsection{Comparación de algoritmos}

\begin{lstlisting}[language=Python,caption={Ajuste y evaluación de los tres modelos.}]
import numpy as np
from sklearn.cluster import KMeans, DBSCAN, AgglomerativeClustering
from sklearn.metrics import silhouette_score, adjusted_rand_score

models = {
    "K-Means":       KMeans(n_clusters=3, n_init=10,
                            random_state=RANDOM_STATE),
    "DBSCAN":        DBSCAN(eps=0.6, min_samples=5),
    "Agglomerative": AgglomerativeClustering(n_clusters=3,
                                             linkage="ward"),
}

for name, m in models.items():
    labels = m.fit_predict(X)
    mask = labels != -1
    sil = (silhouette_score(X[mask], labels[mask])
           if len(set(labels[mask])) > 1 else np.nan)
    ari = adjusted_rand_score(y_true, labels)
    n_noise = int((labels == -1).sum())
    print(f"{name:15s}  sil={sil:.3f}  ari={ari:.3f}  ruido={n_noise}")
\end{lstlisting}
